\section{Agentic Judge Architecture}
\label{app:agentic_judge}

The evaluation system is not a single LLM call but a multi-component
\emph{agentic judge} comprising four dimensional sub-judges, a chronicler
for argument state tracking, a forensics module for stage attribution,
and a meta-judge for bias detection and calibration. This appendix
provides the full architecture.

\subsection{Dimensional Sub-Judges}

Each speech is evaluated by four specialized sub-judges, each with its
own rubric and GRPO weight. Sub-judges operate independently on the
same speech and debate context, then their scores are aggregated.

\begin{table}[h]
\centering
\small
\begin{tabular}{lccp{5.5cm}}
\toprule
\textbf{Sub-Judge} & \textbf{Weight} & \textbf{Scale} & \textbf{Rubric Focus} \\
\midrule
ArgumentJudge & $1.5\times$ & 0--1 & Claim-warrant linkage, logical validity, argument sophistication, burden of proof allocation \\
EvidenceJudge & $1.5\times$ & 0--1 & Citation accuracy (verified against source), source credibility, evidence-argument integration \\
ClashJudge & $1.3\times$ & 0--1 & Response quality to opponent arguments, identification of dropped arguments, attack effectiveness \\
LanguageJudge & $1.0\times$ & 0--1 & Clarity of expression, persuasive impact, signposting and organization, word economy \\
\bottomrule
\end{tabular}
\caption{Dimensional sub-judges with GRPO weights and rubric focus areas.
ArgumentJudge and EvidenceJudge receive higher weight ($1.5\times$) because
argument quality and evidence integrity are most predictive of debate outcomes.}
\label{tab:sub_judges}
\end{table}

The composite inline score for speech $s$ is:
\begin{equation}
\text{score}_{\text{inline}}(s) = \frac{\sum_{j} w_j \cdot \text{score}_j(s)}{\sum_{j} w_j}
\end{equation}

\noindent where $j$ indexes the four sub-judges and $w_j$ is the GRPO weight.

\subsection{Speech-Specific Primary Dimensions}

Beyond the four universal dimensions, each IPDA speech type has a
\emph{primary dimension} reflecting its unique role in the debate structure.
The primary dimension receives additional weight in the composite score
for that speech type.

\begin{table}[h]
\centering
\small
\begin{tabular}{llcp{5cm}}
\toprule
\textbf{Speech} & \textbf{Primary Dimension} & \textbf{Add'l Weight} & \textbf{Evaluation Criteria} \\
\midrule
AC & \texttt{prima\_facie\_case} & 0.20 & Establishes complete burden structure with independent contentions, clear link chain, and sufficient warranting \\
NC & \texttt{offense\_defense\_balance} & 0.18 & Strategic time allocation between attacking the AC and building independent negative offense \\
1AR & \texttt{triage\_efficiency} & 0.18 & Prioritized response to NC arguments; addresses highest-impact attacks first within time constraints \\
NR & \texttt{line\_picture\_synthesis} & 0.16 & Extends negative arguments, synthesizes the round into a coherent ``big picture'' narrative \\
2AR & \texttt{crystallization\_cascade} & 0.18 & Reduces the entire debate to 2--3 decisive voting issues with clear impact calculus \\
\midrule
AC\_CX\_Q & \texttt{case\_probe\_extraction} & 0.15 & Negative questions probing the fresh AC for definitional weaknesses and unstated assumptions \\
AC\_CX\_A & \texttt{case\_fortification} & 0.15 & Affirmative answers that protect case structure without conceding strategic ground \\
NC\_CX\_Q & \texttt{offense\_limitation} & 0.15 & Affirmative questions limiting the scope and impact of negative attacks \\
NC\_CX\_A & \texttt{offense\_preservation} & 0.15 & Negative answers defending attacks while avoiding admissions that weaken offense \\
\bottomrule
\end{tabular}
\caption{Speech-specific primary dimensions in the IPDA format.
Constructive speeches (AC, NC) and final rebuttals (2AR) receive higher
additional weight because they have the strongest structural requirements.}
\label{tab:speech_dimensions}
\end{table}

\subsection{Chronicler: Argument State Tracking}

The Chronicler module maintains a running record of every argument's
status throughout the debate, enabling the ClashJudge to accurately
identify dropped arguments and the BackpropScorer
(Appendix~\ref{app:trajectory}) to attribute outcomes to specific
decisions.

Each argument is tracked through six possible states:

\begin{table}[h]
\centering
\small
\begin{tabular}{lp{8cm}}
\toprule
\textbf{State} & \textbf{Definition} \\
\midrule
\textsc{Standing} & Introduced and not yet addressed by the opponent \\
\textsc{Attacked} & Directly refuted or challenged by the opponent \\
\textsc{Dropped} & Not addressed in the opponent's response when it should have been; typically conceded by omission \\
\textsc{Extended} & Reaffirmed and developed further by the originating side \\
\textsc{Conceded} & Explicitly acknowledged as valid by the opponent \\
\textsc{Turned} & Opponent has reframed the argument to support their own position \\
\bottomrule
\end{tabular}
\caption{Argument state taxonomy maintained by the Chronicler.}
\label{tab:argument_states}
\end{table}

The Chronicler processes speeches sequentially, updating its argument
ledger after each speech. This ledger is provided as context to all
sub-judges and to the BackpropScorer, ensuring that evaluations reflect
the actual flow of the debate rather than assessing each speech in isolation.

\subsection{Forensics: Stage Attribution}

When a speech receives a low score, the Forensics module identifies
\emph{which pipeline stage} caused the failure. This decomposition is
critical for group-wise GRPO (Appendix~\ref{app:grpo_groups}), as it
determines which training group should receive the penalty signal.

\begin{table}[h]
\centering
\small
\begin{tabular}{lp{6cm}c}
\toprule
\textbf{Stage} & \textbf{Failure Indicators} & \textbf{GRPO Group} \\
\midrule
Tactic selection & Wrong strategy for debate state; misread opponent's position & Group A \\
Skeleton building & Poor contention structure; missing link in argument chain & Group B \\
Evidence selection & Irrelevant evidence; fabricated citations; misquoted sources & Group B \\
Speech execution & Good plan but poor delivery; verbose or unclear prose & Group C \\
\bottomrule
\end{tabular}
\caption{Forensics stage attribution mapping failures to GRPO training groups.}
\label{tab:forensics}
\end{table}

The Forensics module produces decomposed rewards:
$r_{\text{tactic}}$, $r_{\text{skeleton}}$, $r_{\text{evidence}}$,
and $r_{\text{execution}}$. These are routed to the appropriate training
group, enabling targeted optimization. For example, a speech with excellent
argumentation but fabricated evidence receives a high $r_{\text{execution}}$
but a strongly negative $r_{\text{evidence}}$, penalizing Group~B without
degrading Group~C's training signal.

\subsection{MetaJudge: Bias Detection and Calibration}

LLM judges exhibit systematic biases that, if uncorrected, contaminate
training data. The MetaJudge detects and mitigates six bias types:

\begin{table}[h]
\centering
\small
\begin{tabular}{lp{7cm}}
\toprule
\textbf{Bias Type} & \textbf{Description} \\
\midrule
\textsc{No\_Alternative\_Penalty} & Penalizing an argument without considering the opponent's alternative \\
\textsc{Extension\_Error} & Treating a re-stated argument as a new argument (inflating or deflating scores) \\
\textsc{Evidence\_Asymmetry} & Applying different evidentiary standards to the two sides \\
\textsc{Recency\_Bias} & Overweighting later speeches relative to earlier ones \\
\textsc{Verbosity\_Bias} & Rewarding longer speeches independent of content quality \\
\textsc{Side\_Bias} & Systematic preference for affirmative or negative regardless of performance \\
\bottomrule
\end{tabular}
\caption{Six bias types detected by the MetaJudge.}
\label{tab:bias_types}
\end{table}

\paragraph{UniversalBiasTracker.}
The MetaJudge maintains a rolling calibration window of the most recent
50 debates. For each judge model and each bias type, it computes the
empirical bias rate and applies score corrections when the rate exceeds
a threshold ($> 0.15$ for most bias types). The calibration window is
updated after each scored debate, enabling the system to adapt to
distributional shift as the trained model improves across iterations.

\paragraph{Multi-model retrospective panel.}
For tournament evaluation and high-stakes training data curation, a
3-model retrospective panel (Claude Opus, Gemini, GPT-4) independently
evaluates each debate. Points of agreement across all three models
(``agreed key moments'') are flagged as the highest-confidence training
signal and receive elevated weight in GRPO.

\subsection{Judge-the-Judge Rubric}

To validate judge quality, a meta-evaluation rubric scores each judge
decision across five dimensions:

\begin{table}[h]
\centering
\small
\begin{tabular}{lcp{6cm}}
\toprule
\textbf{Dimension} & \textbf{Weight} & \textbf{Criteria} \\
\midrule
Evidentiary fidelity & 25\% & Does the judge accurately reference what was actually said in speeches? \\
Symmetric treatment & 20\% & Are the same standards applied to both sides? \\
Clash resolution & 25\% & Does the judge correctly identify who won contested arguments? \\
Decision clarity & 15\% & Is the decision logically derived from the analysis? \\
Technical accuracy & 15\% & Does the judge correctly apply debate norms (dropped arguments, burden of proof, etc.)? \\
\bottomrule
\end{tabular}
\caption{Judge-the-judge rubric for meta-evaluation of scoring quality.}
\label{tab:judge_rubric}
\end{table}

Judges scoring below 0.60 on the meta-rubric have their scores discarded
for that debate. In tournament evaluation, this filtering removed
approximately 8\% of individual judge decisions, predominantly due to
evidentiary fidelity failures (the judge misattributed or fabricated
quotes from speeches).
